%-----------------------------------------------------------------------
\section{Discussion}
\label{sec:discussion}
%-----------------------------------------------------------------------

\subsection{What the comparison actually established}
\label{sec:disc-what}

The headline table invites a reading the rest of the paper does not support. On
the coarsest, most stable scale our field correlates better with an independent
trace than either reference solver does, and at matched splat budgets above
roughly \num{1000000} it is both cheaper and closer than plain geometric
sampling. Neither of those is a general speed-up claim. Below that budget the
geometric estimator wins outright, as
Table~\ref{tab:efficiency} shows, and on the finest scale measured SPPM ties us.

The defensible statement is narrower and, we think, more interesting. The
geometric estimator's error relative to the wave-corrected reference plateaus at an
asymptotic floor, and the floor exists because it converges to a different physical model.
Rays carry zero energy into wave diffraction fringes, so no geometric sampling budget
reaches the wave-regularized target. Refining a sampler converges the estimate of the
geometric integral and cannot change which physical model is being evaluated.
That is a statement about cross-model reachability rather than about solver speed,
and it is the one the data supports.

\subsection{The estimator keeps being the problem}
\label{sec:disc-estimator}

The three results we withdrew or corrected from earlier drafts were all killed by
estimators rather than by core physics: a correlation that was an artefact of a
sign convention and a tile-naming inversion, a dispersion model that had been fitted
rather than measured, and a flux-conservation figure that was circular by construction.
All are reported in full because they share a structure that a reader
reproducing this work will meet.

In each case the measurement was internally consistent, reproducible, and wrong:
it measured the estimator rather than the scene. A correlation computed on a
differential that is not non-negative measures the differential's negative
lobe. An exponent fitted over three wavelengths measures the baseline. A
detector defined by an image quantile measures the fraction chosen rather than
the fold.

We do not think this is a peculiarity of caustics. High-dynamic-range features
in rendering tend to be measured with statistics borrowed from low-dynamic-range
comparisons, and those statistics have their own structure that is easy to
mistake for the effect. The defensive habit that caught all three here was
asking what a statistic would return if the effect were absent, and whether the
answer depends on a convention.

\subsection{Guidance for a practitioner}
\label{sec:disc-guidance}

Three things from this work transfer without the rest.

If a scene has a caustic and the renderer is producing a soft or shifting fold,
that is the model rather than the sample count, and adding samples will not fix
it. Raising the budget makes the softness converge toward a wrong shape.

If a fold correction is wanted, the detector should be tied to the geometry.
Thresholding $|\detJ|$ converges under refinement; thresholding image brightness
does not, and the failure is invisible in a single frame because the flagged
region looks plausible either way. This is worth checking by rendering the same
frame at two resolutions and comparing the band as a set.

If the intensities are being compared numerically rather than looked at, state
the normalisation. Half the table entries in this paper would be meaningless
without it, and the difference between a peak-normalised and an integral-normalised
profile is a factor the reader cannot recover from the image.

\subsection{The coherence question, restated}
\label{sec:disc-coherence}

The limitation of Section~\ref{sec:limits-coherence} deserves repeating in the
discussion because it changes how the paper should be read. The fold profile
computed here is the coherent-limit profile. Under a source of the size used in
this scene, the physical fringe spacing ($d \approx 6.8\,\mu\text{m}$) is roughly \num{136} times smaller than the
source's geometric blur ($925\,\mu\text{m}$), so a photograph would not show those fringes.

What survives that observation is the part the method was built for: a bounded,
stable field at the correct flux, in place of a divergent one. The fringes are
the visible signature of the correction but not its justification, and a reader
interested in the justification should read
Table~\ref{tab:efficiency} and Table~\ref{tab:flicker} rather than
Figure~\ref{fig:airy}. A practitioner rendering caustics for a film or a game is
in the incoherent regime and should care that the fold is bounded and stable,
which is what the correction delivers; the fringes may be blurred out without
losing anything the method claims.

\subsection{Cost, and where it goes}
\label{sec:disc-cost}

Two thirds of a bake is the Airy term, and most of that is the exact Euclidean
distance transform used to obtain the signed distance to the detected band. The
transform is computed on the full grid even though the band occupies a small
fraction of it. Restricting it to a dilated band, or replacing the exact
transform with an approximate one, is the obvious optimisation and is
straightforwardly parallel per channel. We have not done it, so the figures in
Table~\ref{tab:cost} are the unoptimised ones, and anyone extending this
implementation should expect to make it substantially faster without touching
the parts of the method that matter.

The other cost is scene-dependent in a way worth flagging. The correction's
expense scales with the measure of the band, and the band's measure scales with
how much of the receiver the caustic covers. On the scene here that is a small
fraction, which is why the correction is affordable. A scene whose caustic
covers a large part of the receiver would change the trade, and we have not
measured where the crossover lies (Section~\ref{sec:limits-cost}).

\subsection{The shape of a follow-up}
\label{sec:disc-followup}

Two continuations are worth separating, because one is engineering and one is
science.

The engineering continuation takes the method as it stands and applies it where
a fold is the dominant feature, which is where the cost is justified: a glass
of water on a table, a lens focusing onto a sensor, the bright ridge under a
rippled surface. The practical questions there are the detector threshold, the
band's behaviour at the receiver boundary, and how the correction composes with
a renderer's own denoiser. All three are measurable.

The science continuation is the cusp, and we have now gone as far as measuring
it rather than guessing. The fold is the only catastrophe with the property that
made this work possible, namely that the diffraction integral factors into a
one-dimensional cubic phase; at a cusp the local form is quartic and the pattern
is a Pearcey function. We built that case and tested the recipe on it
(Section~\ref{sec:verify-cusp}). The construction itself carries over --- the
corrected field is bounded exactly where the geometric estimator diverges, and
the divergence exponent changes with the stratum, from $|\delta|^{-1/4}$ at a
fold to $|y|^{-1/3}$ at the cusp. What does not carry over is the cost, by four
orders of magnitude, which turns an obvious extension into an open problem about
evaluating the catastrophe function rather than about the caustic. That is a
narrower claim than we would like and a more useful one than leaving it open.
